\section{P04 - Corrigé - Types construits complément}

\subsection{Capacités officielles atomiques}

\begin{itemize}
\tightlist
\item
  P-DATA-CONSTR-01 : Écrire une fonction renvoyant un p-uplet de valeurs
\item
  P-DATA-CONSTR-02B : Construire un tableau par compréhension
\item
  P-DATA-CONSTR-02C : Utiliser des tableaux de tableaux pour représenter
  des matrices
\item
  P-DATA-CONSTR-03A : Construire une entrée de dictionnaire
\item
  P-DATA-CONSTR-03B : Itérer sur les éléments d'un dictionnaire
\end{itemize}

\begin{center}\rule{0.5\linewidth}{0.5pt}\end{center}

\subsection{Corrigé du TD}

\subsubsection{Exercice 1 - Fonction renvoyant un p-uplet (P-DATA-CONSTR-01)}

\textbf{1a.} La valeur renvoyée est de type \texttt{tuple}. Plus
précisément, c'est un tuple de trois éléments :
\texttt{(8,\ 15,\ 11.25)}.

\textbf{1b.}

\begin{itemize}
\tightlist
\item
  \texttt{a\ =\ 8} (minimum)
\item
  \texttt{b\ =\ 15} (maximum)
\item
  \texttt{c\ =\ 11.25} (moyenne : (12 + 8 + 15 + 10) / 4 = 45 / 4 =
  11.25)
\end{itemize}

\textbf{1c.}

\begin{python}
def extremes(a, b):
    if a <= b:
        return a, b
    else:
        return b, a
\end{python}

\textbf{1d.} L'appel \texttt{analyser\_notes({[}{]})} provoque une
erreur \texttt{IndexError} sur la ligne \texttt{mini\ =\ notes{[}0{]}}
car le tableau est vide et l'indice 0 n'existe pas. De plus,
\texttt{len(notes)} vaut 0, ce qui provoquerait une division par zéro.
La précondition est : le tableau doit être non vide.

\begin{center}\rule{0.5\linewidth}{0.5pt}\end{center}

\subsubsection{Exercice 2 - Tableau par compréhension (P-DATA-CONSTR-02B)}

\textbf{2a.}

\begin{itemize}
\tightlist
\item
  \texttt{L1\ =\ {[}0,\ 1,\ 4,\ 9,\ 16,\ 25,\ 36,\ 49,\ 64,\ 81{]}}
\item
  \texttt{L2\ =\ {[}3,\ 6,\ 9,\ 12,\ 15{]}}
\item
  \texttt{L3\ =\ {[}\textquotesingle{}P\textquotesingle{},\ \textquotesingle{}Y\textquotesingle{},\ \textquotesingle{}T\textquotesingle{},\ \textquotesingle{}H\textquotesingle{},\ \textquotesingle{}O\textquotesingle{},\ \textquotesingle{}N\textquotesingle{}{]}}
\end{itemize}

\textbf{2b.}

\texttt{L4\ =\ {[}3,\ 6,\ 9,\ 12,\ 15,\ 18{]}} (les multiples de 3 entre
1 et 20).

\textbf{2c.}

\begin{python}
cubes = [x ** 3 for x in range(1, 9)]
\end{python}

\textbf{2d.}

\begin{python}
positives = [t for t in temperatures if t > 0]
# Résultat : [15.5, 22.0, 18.3, 30.2]
\end{python}

\textbf{2e.} \texttt{{[}0\ for\ \_\ in\ range(0){]}} produit
\texttt{{[}{]}} (liste vide). \texttt{range(0)} génère une séquence
vide, donc la compréhension n'itère sur aucun élément.

\begin{center}\rule{0.5\linewidth}{0.5pt}\end{center}

\subsubsection{Exercice 3 - Matrices (P-DATA-CONSTR-02C)}

\textbf{3a.}

\begin{itemize}
\tightlist
\item
  \texttt{M{[}0{]}{[}2{]}\ =\ 8} (ligne 0, colonne 2)
\item
  \texttt{M{[}1{]}{[}1{]}\ =\ 7} (ligne 1, colonne 1)
\item
  \texttt{M{[}2{]}{[}0{]}\ =\ 9} (ligne 2, colonne 0)
\end{itemize}

\textbf{3b.} \texttt{M} possède 3 lignes et 3 colonnes.

\begin{itemize}
\tightlist
\item
  Nombre de lignes : \texttt{len(M)} qui vaut \texttt{3}.
\item
  Nombre de colonnes : \texttt{len(M{[}0{]})} qui vaut \texttt{3}.
\end{itemize}

\textbf{3c.} Trace d'exécution de \texttt{somme\_diagonale(M)} :

\begin{longtable}[]{@{}llll@{}}
\toprule\noalign{}
Étape & \texttt{i} & \texttt{matrice{[}i{]}{[}i{]}} & \texttt{total} \\
\midrule\noalign{}
\endhead
\bottomrule\noalign{}
\endlastfoot
Init & - & - & 0 \\
1 & 0 & 5 & 5 \\
2 & 1 & 7 & 12 \\
3 & 2 & 6 & 18 \\
\end{longtable}

Résultat renvoyé : \texttt{18}.

\textbf{3d.} Le code \texttt{grille\ =\ {[}{[}0{]}\ *\ 3{]}\ *\ 3} crée
une liste contenant 3 fois la \textbf{même} sous-liste (même objet en
mémoire). Modifier \texttt{grille{[}0{]}{[}0{]}} modifie donc les trois
lignes simultanément.

Affichage :
\texttt{{[}{[}1,\ 0,\ 0{]},\ {[}1,\ 0,\ 0{]},\ {[}1,\ 0,\ 0{]}{]}}.

Correction :

\begin{python}
grille = [[0 for j in range(3)] for i in range(3)]
grille[0][0] = 1
# Résultat correct : [[1, 0, 0], [0, 0, 0], [0, 0, 0]]
\end{python}

\begin{center}\rule{0.5\linewidth}{0.5pt}\end{center}

\subsubsection{Exercice 4 - Construire une entrée de dictionnaire (P-DATA-CONSTR-03A)}

\textbf{4a.}

\begin{python}
photo = {}
photo["marque"] = "Sony"
photo["modele"] = "A7III"
photo["date_prise"] = "2025-07-01"
photo["taille"] = (6000, 4000)
\end{python}

\textbf{4b.} \texttt{photo{[}"taille"{]}} affiche
\texttt{(6000,\ 4000)}. C'est un \texttt{tuple} de deux entiers.

\textbf{4c.} Après \texttt{photo{[}"iso"{]}\ =\ 100} puis
\texttt{photo{[}"iso"{]}\ =\ 200}, la clé \texttt{"iso"} est associée à
\texttt{200}. La seconde affectation remplace la valeur précédente. Un
dictionnaire ne peut pas contenir deux entrées avec la même clé.

\textbf{4d.} L'instruction \texttt{d{[}{[}1,\ 2{]}{]}\ =\ "coordonnees"}
provoque
\texttt{TypeError:\ unhashable\ type:\ \textquotesingle{}list\textquotesingle{}}.
Les clés d'un dictionnaire doivent être de type immuable (hashable). Une
liste est mutable et ne peut donc pas servir de clé. On pourrait
utiliser un tuple a la place :
\texttt{d{[}(1,\ 2){]}\ =\ "coordonnees"}.

\begin{center}\rule{0.5\linewidth}{0.5pt}\end{center}

\subsubsection{Exercice 5 - Itérer sur un dictionnaire (P-DATA-CONSTR-03B)}

\textbf{5a.}

Bloc A affiche (une par ligne) :

\begin{python}
temperature
humidite
vent
pression
\end{python}

Bloc B affiche :

\begin{python}
22
65
15
1013
\end{python}

Bloc C affiche :

\begin{python}
temperature = 22
humidite = 65
vent = 15
pression = 1013
\end{python}

\textbf{5b.}

\begin{python}
total = 0
for val in meteo.values():
    total = total + val
# total vaut 22 + 65 + 15 + 1013 = 1115
\end{python}

\textbf{5c.}

\begin{python}
def cles_superieures(d, seuil):
    resultat = []
    for cle, valeur in d.items():
        if valeur > seuil:
            resultat.append(cle)
    return resultat
\end{python}

Exemple : \texttt{cles\_superieures(meteo,\ 20)} renvoie
\texttt{{[}"temperature",\ "humidite",\ "pression"{]}}.

\textbf{5d.} Itérer sur un dictionnaire vide \texttt{\{\}} ne provoque
pas d'erreur. La boucle \texttt{for} ne s'exécute simplement pas car il
n'y a aucun élément à parcourir.

\begin{center}\rule{0.5\linewidth}{0.5pt}\end{center}

\subsection{Corrigé de l’évaluation}

\subsubsection{Question 1 - P-DATA-CONSTR-01 (4 points)}

\textbf{1a.} (1 pt) La valeur renvoyée par \texttt{encadrer(5)} est de
type \texttt{tuple}. C'est le tuple \texttt{(4,\ 6)}.

\textbf{1b.} (1 pt) \texttt{a\ =\ 9}, \texttt{b\ =\ 11}.

\textbf{1c.} (2 pts)

\begin{python}
def aire_perimetre(longueur, largeur):
    aire = longueur * largeur
    perimetre = 2 * (longueur + largeur)
    return aire, perimetre
\end{python}

Pour \texttt{longueur\ =\ 5} et \texttt{largeur\ =\ 3} :
\texttt{aire\_perimetre(5,\ 3)} renvoie \texttt{(15,\ 16)}.

\begin{center}\rule{0.5\linewidth}{0.5pt}\end{center}

\subsubsection{Question 2 - P-DATA-CONSTR-02B (4 points)}

\textbf{2a.} (2 pts)

\begin{itemize}
\tightlist
\item
  \texttt{L1\ =\ {[}0,\ 1,\ 4,\ 9,\ 16,\ 25,\ 36,\ 49,\ 64,\ 81{]}}
\item
  \texttt{L2\ =\ {[}5,\ 10,\ 15{]}}
\end{itemize}

\textbf{2b.} (2 pts)

\begin{python}
notes = [8, 12, 5, 16, 10, 3]
reussies = [n for n in notes if n >= 10]
# Résultat : [12, 16, 10]
\end{python}

\begin{center}\rule{0.5\linewidth}{0.5pt}\end{center}

\subsubsection{Question 3 - P-DATA-CONSTR-02C (4 points)}

\textbf{3a.} (1 pt) \texttt{grille{[}1{]}{[}2{]}\ =\ 5},
\texttt{grille{[}2{]}{[}0{]}\ =\ 4}.

\textbf{3b.} (1 pt) Nombre de lignes : \texttt{len(grille)} (vaut 3).
Nombre de colonnes : \texttt{len(grille{[}0{]})} (vaut 3).

\textbf{3c.} (2 pts)

\begin{python}
def somme_ligne(matrice, i):
    total = 0
    for j in range(len(matrice[i])):
        total = total + matrice[i][j]
    return total
\end{python}

Pour la ligne 0 de \texttt{grille} : \texttt{somme\_ligne(grille,\ 0)}
renvoie \texttt{2\ +\ 7\ +\ 1\ =\ 10}.

\begin{center}\rule{0.5\linewidth}{0.5pt}\end{center}

\subsubsection{Question 4 - P-DATA-CONSTR-03A (4 points)}

\textbf{4a.} (2 pts)

\begin{python}
eleve = {}
eleve["nom"] = "Martin"
eleve["prenom"] = "Léa"
eleve["age"] = 16
eleve["classe"] = "1NSI"
\end{python}

\textbf{4b.} (1 pt) L'instruction \texttt{eleve{[}"age"{]}\ =\ 17}
remplace la valeur 16 par 17. La nouvelle valeur de
\texttt{eleve{[}"age"{]}} est \texttt{17}.

\textbf{4c.} (1 pt) Une liste est un type mutable et ne peut pas être
utilisée comme clé de dictionnaire. Python exige des clés hashables
(immuables). L'erreur levée est
\texttt{TypeError:\ unhashable\ type:\ \textquotesingle{}list\textquotesingle{}}.

\begin{center}\rule{0.5\linewidth}{0.5pt}\end{center}

\subsubsection{Question 5 - P-DATA-CONSTR-03B (4 points)}

\textbf{5a.} (1 pt) Le code affiche :

\begin{python}
Alice
Bob
Clara
\end{python}

\textbf{5b.} (1 pt)

\begin{python}
for nom, score in scores.items():
    print(f"{nom} a obtenu {score}")
\end{python}

\textbf{5c.} (2 pts)

\begin{python}
def meilleur_score(d):
    meilleur_nom = None
    meilleur_val = None
    for nom, score in d.items():
        if meilleur_val is None or score > meilleur_val:
            meilleur_nom = nom
            meilleur_val = score
    return meilleur_nom, meilleur_val
\end{python}

Avec \texttt{scores} : \texttt{meilleur\_score(scores)} renvoie
\texttt{("Clara",\ 91)}.

\begin{center}\rule{0.5\linewidth}{0.5pt}\end{center}

\subsection{Corrigé du TP}

\subsubsection{Exercice 1}

\textbf{1a.}

\begin{python}
def statistiques(tableau):
    mini = tableau[0]
    maxi = tableau[0]
    somme = 0
    for val in tableau:
        if val < mini:
            mini = val
        if val > maxi:
            maxi = val
        somme = somme + val
    moyenne = somme / len(tableau)
    return mini, maxi, moyenne
\end{python}

\textbf{1b.}

\begin{python}
def secondes_vers_hms(total_secondes):
    heures = total_secondes // 3600
    reste = total_secondes % 3600
    minutes = reste // 60
    secondes = reste % 60
    return heures, minutes, secondes
\end{python}

\textbf{1c.} \texttt{statistiques({[}{]})} provoque \texttt{IndexError}
car \texttt{tableau{[}0{]}} est inaccessible sur un tableau vide.

\begin{center}\rule{0.5\linewidth}{0.5pt}\end{center}

\subsubsection{Exercice 2}

\textbf{2a.}

\begin{python}
carres = [x ** 2 for x in range(10)]
multiples_7 = [k for k in range(0, 51) if k % 7 == 0]
longueurs = [len(m) for m in mots]
\end{python}

\textbf{2b.}

\begin{python}
pairs_carres = [n ** 2 for n in nombres if n % 2 == 0]
\end{python}

\textbf{2c.}

\begin{python}
def initiales(prenoms):
    return [p[0].upper() for p in prenoms]
\end{python}

\begin{center}\rule{0.5\linewidth}{0.5pt}\end{center}

\subsubsection{Exercice 3}

\textbf{3a.}

\begin{python}
identite = [[1 if i == j else 0 for j in range(3)] for i in range(3)]
\end{python}

\textbf{3b.}

\begin{python}
def transposer(matrice):
    nb_lignes = len(matrice)
    nb_colonnes = len(matrice[0])
    resultat = [[0 for j in range(nb_lignes)] for i in range(nb_colonnes)]
    for i in range(nb_lignes):
        for j in range(nb_colonnes):
            resultat[j][i] = matrice[i][j]
    return resultat
\end{python}

\textbf{3c.}

\begin{python}
def rechercher(matrice, valeur):
    for i in range(len(matrice)):
        for j in range(len(matrice[i])):
            if matrice[i][j] == valeur:
                return (i, j)
    return None
\end{python}

\begin{center}\rule{0.5\linewidth}{0.5pt}\end{center}

\subsubsection{Exercice 4}

\textbf{4a.}

\begin{python}
film = {}
film["titre"] = "Interstellar"
film["realisateur"] = "Christopher Nolan"
film["annee"] = 2014
film["duree_min"] = 169
film["genres"] = ["science-fiction", "drame"]
\end{python}

\textbf{4b.}

\begin{python}
def compter_mots(phrase):
    mots = phrase.split()
    compteur = {}
    for mot in mots:
        if mot in compteur:
            compteur[mot] = compteur[mot] + 1
        else:
            compteur[mot] = 1
    return compteur
\end{python}

\textbf{4c.}

\begin{python}
def creer_exif(appareil, date, largeur, hauteur, iso):
    exif = {}
    exif["appareil"] = appareil
    exif["date"] = date
    exif["resolution"] = (largeur, hauteur)
    exif["iso"] = iso
    return exif
\end{python}

\begin{center}\rule{0.5\linewidth}{0.5pt}\end{center}

\subsubsection{Exercice 5}

\textbf{5a.}

\begin{python}
def afficher_fiche(d):
    for cle, valeur in d.items():
        print(f"{cle} : {valeur}")
\end{python}

\textbf{5b.}

\begin{python}
def filtrer_valeurs(d, seuil):
    resultat = {}
    for cle, valeur in d.items():
        if valeur >= seuil:
            resultat[cle] = valeur
    return resultat
\end{python}

\textbf{5c.}

\begin{python}
def inverser_dict(d):
    resultat = {}
    for cle, valeur in d.items():
        resultat[valeur] = cle
    return resultat
\end{python}

\subsection{Objectifs}

\subsection{Prérequis}

\subsection{Situation-problème}

Un programme de gestion de notes doit stocker, pour chaque élève, son
nom et ses notes dans différentes matières. Quelle structure de données
choisir et comment organiser les accès ?

\subsection{Activité d’entrée}

Écrire une fonction Python qui reçoit deux nombres et renvoie à la fois
leur somme et leur produit. Tester avec print(resultat{[}0{]},
resultat{[}1{]}).

\subsection{Exercices}

Exercices de manipulation de tuples, listes en compréhension, matrices
et dictionnaires.

\subsection{Erreurs fréquentes}

\begin{itemize}
\tightlist
\item
  EF1 : confondre tuple (immutable) et liste (mutable) lors du retour de
  fonction.
\item
  EF2 : oublier les crochets dans une compréhension de liste.
\item
  EF3 : accéder à une clé inexistante d'un dictionnaire sans utiliser
  get().
\end{itemize}

\subsection{Remédiation}

Exercice de remédiation : écrire une fonction min\_max qui renvoie le
minimum et le maximum d'une liste sous forme de tuple.

\subsection{Différenciation}

\begin{itemize}
\tightlist
\item
  Socle : exercices de base.
\item
  Standard : exercices complets.
\item
  Expert : exercice d'approfondissement.
\end{itemize}

\subsection{Critères de réussite}

\begin{itemize}
\tightlist
\item
  Critère de réussite : la fonction renvoie un tuple vérifiable par
  déstructuration.
\item
  Critère de validation : la compréhension de liste produit le résultat
  attendu.
\item
  Observable : le dictionnaire est itéré avec keys(), values() et
  items() correctement.
\end{itemize}

\subsection{Séance(s) correspondante(s)}

Séance dédiée.

\subsubsection{Corrigé exercice 1}

\textbf{Méthode} : on écrit une fonction qui renvoie un tuple contenant
les deux valeurs calculées. La fonction \texttt{min\_max(lst)} renvoie
\texttt{(1,\ 9)} pour la liste \texttt{{[}3,\ 1,\ 9,\ 5{]}}. Le résultat
vaut \texttt{(min,\ max)} sous forme de tuple Python.

\subsubsection{Corrigé exercice 2}

Réponse détaillée avec justification.

La compréhension donne
\texttt{{[}0,\ 1,\ 4,\ 9,\ 16,\ 25,\ 36,\ 49,\ 64,\ 81{]}} pour les
carrés de 0 à 9.

\subsubsection{Corrigé exercice 3}

Réponse détaillée avec justification.

L'élément \texttt{matrice{[}1{]}{[}2{]}} vaut \texttt{5} et
\texttt{somme\_diagonale} renvoie \texttt{18}.

\subsubsection{Corrigé exercice 4}

Réponse détaillée avec justification.

Le dictionnaire donne
\texttt{\{\textquotesingle{}titre\textquotesingle{}:\ \textquotesingle{}Interstellar\textquotesingle{},\ \textquotesingle{}realisateur\textquotesingle{}:\ \textquotesingle{}Christopher\ Nolan\textquotesingle{},\ \textquotesingle{}annee\textquotesingle{}:\ 2014\}}.

\subsubsection{Corrigé exercice 5}

Réponse détaillée avec justification.

\texttt{d.keys()} renvoie
\texttt{dict\_keys({[}\textquotesingle{}temperature\textquotesingle{},\ \textquotesingle{}humidite\textquotesingle{},\ \textquotesingle{}vent\textquotesingle{},\ \textquotesingle{}pression\textquotesingle{}{]})}
(une vue, pas une liste) pour le dictionnaire météo.

\subsubsection{Corrigé exercice 6}

Réponse détaillée avec justification.

La fonction renvoie \texttt{(8,\ 15,\ 11.25)} pour la liste
\texttt{{[}12,\ 8,\ 15,\ 10{]}}.

\subsubsection{Corrigé exercice 7}

Réponse détaillée avec justification.

La compréhension filtrée donne \texttt{{[}12,\ 16,\ 10{]}} pour les
notes supérieures ou égales à 10.

\subsubsection{Corrigé exercice 8}

Réponse détaillée avec justification.

La fonction \texttt{inverser\_dict} renvoie
\texttt{\{\textquotesingle{}a\textquotesingle{}:\ 1,\ \textquotesingle{}b\textquotesingle{}:\ 2\}}
pour l'entrée
\texttt{\{1:\ \textquotesingle{}a\textquotesingle{},\ 2:\ \textquotesingle{}b\textquotesingle{}\}}.
